% preamble-elsarticle.tex -- typography for the Journal of Computational Algebra (Elsevier)
% build.  Loaded by tools/make-jaca.sh right after \documentclass[preprint,12pt]{elsarticle}.
%
% Changelog (reverse-chronological)
%   2026-10-21  Claude-Code-3c0af632 (jts-pc), at Jan's request ("you might as well use the
%               Elsevier style file").  Created as the "one-file swap" preamble.tex was designed
%               for.  Differences from preamble.tex, and why:
%               (a) No fontspec, geometry or KOMA: elsarticle sets the page, and Elsevier's
%                   Editorial Manager compiles LaTeX with pdflatex, so nothing here may need
%                   xelatex.  make-jaca.sh builds with BOTH engines and fails on either.
%               (b) No tcolorbox.  The converted body calls every theorem-like environment as
%                   \begin{theorem}{TITLE}{LABEL} (tcolorbox's signature), so the environments
%                   below keep that signature but are plain amsthm theorems: Elsevier typesets
%                   its own look, and coloured boxes in a submission only get in the way.
%               (c) elsarticle's own thebibliography (natbib, numbered), which prints its own
%                   "References" heading; make-jaca.sh removes the body's \section{References}
%                   instead of redefining the environment as preamble.tex has to.
%               (d) \emergencystretch 6em and 9pt verbatim, for the narrower measure (both
%                   explained where they are set).
%               Everything else -- \thmref, \fitdisplay, the notation macros, the breakable
%               underscore -- is copied from preamble.tex unchanged.  Keep the two in step.

% ---- packages ----------------------------------------------------------
% graphicx is loaded by elsarticle itself.
\usepackage{amsmath,amssymb,amsthm,mathtools}
\usepackage{microtype}
\usepackage[hidelinks]{hyperref}
\usepackage[capitalise,noabbrev]{cleveref}
\usepackage{booktabs}
% pandoc emits longtable for every table, plus these helpers (see preamble.tex).
\usepackage{longtable,array,calc}
\usepackage{etoolbox}
\usepackage{orcidlink}
\usepackage{tikz}
\usetikzlibrary{arrows.meta,positioning,shapes.geometric}
% Converted prose is dense with unbreakable inline maths (see preamble.tex).  6em here, not
% preamble.tex's 3em: elsarticle's measure is 390pt, and at 3em xelatex (Latin Modern metrics)
% left two lines 11pt and 24pt over -- a long inline formula, and a typewriter filename that
% has no break point.  6em clears both under both engines, at the price of two or four loose
% (badness 10000) lines; Elsevier re-typesets the article in any case.
\setlength{\emergencystretch}{6em}
\providecommand{\tightlist}{\setlength{\itemsep}{0pt}\setlength{\parskip}{0pt}}
\providecommand{\pandocbounded}[1]{#1}

% ---- words never to hyphenate ---------------------------------------------
% TeX broke the DOI 10.5281/zenodo.23433729 as "zen-odo" at a line end, and a hyphen inside a
% DOI gets copied as part of it.  An exception with no break points forbids it.
\hyphenation{zenodo}

% ---- verbatim -------------------------------------------------------------
% The appendix reproduces raw program output in verbatim, with lines up to ~76 characters.
% At this class's 12pt and 390pt measure the widest ran 79pt into the margin; 9pt cmtt fits
% them all (cmtt has a real 9pt design size, so nothing is scaled or substituted).
\makeatletter
\renewcommand{\verbatim@font}{\fontsize{9}{11}\selectfont\ttfamily}
\makeatother

% ---- theorem-like environments -----------------------------------------
% ONE shared counter for everything except Question, exactly as in preamble.tex and in the
% Typst source: the numbers (Theorem 6, Remark 45, ...) are cited by name in the Lean sources
% and elsewhere, so the three routes must agree.  Question is deliberately off the counter.
\theoremstyle{plain}
\newtheorem{jtsTheorem}{Theorem}
\newtheorem{jtsProposition}[jtsTheorem]{Proposition}
\newtheorem{jtsLemma}[jtsTheorem]{Lemma}
\theoremstyle{definition}
\newtheorem{jtsDefinition}[jtsTheorem]{Definition}
\newtheorem{jtsExample}[jtsTheorem]{Example}
\newtheorem{jtsProblem}[jtsTheorem]{Problem}
\newtheorem{jtsQuestion}{Question}
\theoremstyle{remark}
\newtheorem{jtsRemark}[jtsTheorem]{Remark}
% \begin{theorem}{TITLE}{LABEL}: an empty TITLE prints no parenthesis; the label gets the
% ":" prefix that tcolorbox prepends, because \thmref below (and every cross-reference that
% tools/markers2env.py writes) looks for ":LABEL".  The title is braced inside [...] so that
% a "]" in it cannot end the optional argument early.
\newcommand{\jtsBox}[2]{%
  \newenvironment{#1}[2]%
    {\ifstrempty{##1}{\begin{#2}}{\begin{#2}[{##1}]}\label{:##2}}%
    {\end{#2}}}
\jtsBox{theorem}{jtsTheorem}
\jtsBox{proposition}{jtsProposition}
\jtsBox{lemma}{jtsLemma}
\jtsBox{definition}{jtsDefinition}
\jtsBox{example}{jtsExample}
\jtsBox{problem}{jtsProblem}
\jtsBox{question}{jtsQuestion}
\jtsBox{remark}{jtsRemark}
\renewcommand{\qedsymbol}{$\square$}

% ---- cross-references into the appendix --------------------------------
% elsarticle's \appendix makes the section number itself read "Appendix A" ("Appendix A.2"
% for a subsection), so cleveref's default printed "Section Appendix A.2".  After
% \appendix, section and subsection labels get a type whose format is the bare number.
\apptocmd{\appendix}{\crefalias{section}{jtsappendix}\crefalias{subsection}{jtsappendix}}{}{%
  \PackageError{preamble-elsarticle}{could not patch \string\appendix}{}}
\crefformat{jtsappendix}{#2#1#3}
\Crefformat{jtsappendix}{#2#1#3}

% ---- cross-references into the theorem-like environments ---------------
% Copied from preamble.tex, where the reasons are written out: the environment name is passed
% in because all of them share one counter, and there are no square brackets in the call.
\newcommand{\thmref}[2]{\hyperref[:#2]{#1~\ref*{:#2}}}

% ---- displays that are wider than the measure (from preamble.tex) -----
\makeatletter
\newsavebox{\jts@fitbox}
\newcommand{\fitdisplay}[1]{%
  \sbox{\jts@fitbox}{\ensuremath{\displaystyle #1}}%
  \ifdim\wd\jts@fitbox>\linewidth
    \resizebox{\linewidth}{!}{\usebox{\jts@fitbox}}%
  \else
    \usebox{\jts@fitbox}%
  \fi}
\makeatother

% ---- notation (from preamble.tex) --------------------------------------
\newcommand{\FF}{\mathbb{F}}
\newcommand{\QQ}{\mathbb{Q}}
\newcommand{\RR}{\mathbb{R}}
\newcommand{\AAA}{\mathbb{A}}
\newcommand{\PP}{\mathbb{P}}
\newcommand{\vek}[1]{\overline{#1}}
% Breakable underscores, for the script filenames in running prose (see preamble.tex).
\let\jtsOldUnderscore\_
\renewcommand{\_}{\jtsOldUnderscore\allowbreak}
\newcommand{\enc}{\operatorname{enc}}
\newcommand{\red}{\operatorname{red}}
\DeclareMathOperator{\tr}{tr}
\DeclareMathOperator{\pd}{pd}
